% !TeX root = ../main.tex

\section{Simulation Setup and Results}

The evaluation will compare all methods on shared LEO downlink scenario
instances. The planned baselines are nearest-visible-satellite association,
best-channel association without handover control, soft handover-penalized
association, full-horizon mixed-integer association, and rolling-horizon
association. When the proxy is separable across cells, a dynamic-programming
solver can also be included as an exact reference for the surrogate objective.

\begin{table}[t]
\caption{Primary Evaluation Metrics}
\label{tab:metrics}
\centering
\begin{tabular}{ll}
\toprule
Metric & Purpose \\
\midrule
Mean satisfaction & Average served-demand ratio \\
Fifth-percentile satisfaction & Worst-cell service quality \\
Log utility & Proportional fairness \\
Total served demand & System throughput \\
Total handovers & Mobility overhead \\
Budget violations & Hard-constraint feasibility \\
Solve time & Algorithm scalability \\
\bottomrule
\end{tabular}
\end{table}

The main metrics are summarized in Table~\ref{tab:metrics}. The final WCNC
claim should be limited to scenarios where the proposed method has zero
avoidable handover-budget violations. Sensitivity studies should vary:
\begin{itemize}
\item demand scale and demand imbalance across cells;
\item visibility density and minimum elevation threshold;
\item handover budget $H_c$;
\item rolling-window length and commit step;
\item proxy calibration error relative to the inner convex allocator.
\end{itemize}
